\section{Conformal inversion and spatial moments of the cubic representative}
\label{cm:section}

\paragraph{Status and scope.}
This is a complete component derivation, author-checked with independent
mathematical audit pending. It resolves two spatial integrability questions
under an explicit weighted-gradient hypothesis: the minimizing representative
belongs to $L^2$, and its scalar divergence defect belongs to $L^{3/2}$.
A weighted energy estimate then supplies uniform finite-horizon spacetime
bounds on every original Schwartz-data classical branch, even below a
putative finite maximal time. It supplies no critical time-integrability
estimate and does not promote HIGH-PRESSURE, HIGH-STRAIN, or DEFECT-L4.

The mechanism is conformal invariance at exponent equal to dimension,
not a general weighted Calder\'on--Zygmund theorem. Conformal invariance
of $n$-harmonic equations is classical; see Liimatainen--Salo
\cite{LiimatainenSalo2014}, introduction, and the nonlinear Hodge context
of~\cite{IwaniecScottStroffolini1999}. Our representative generally has
nonzero curl, so no scalar $n$-harmonic regularity theorem is applied.
All transformation, density, and moment arguments needed here are proved
below. In particular, Navier--Stokes is not asserted to be invariant under
this spatial inversion.

\subsection{An isometry of the gradient cosets}

For $x\ne0$ put
\[
 I(x)=\frac{x}{|x|^2},\qquad n_x=\frac{x}{|x|},\qquad
 R_x=\operatorname{Id}-2n_x\otimes n_x,\qquad
 (\mathcal Kf)(x)=|x|^{-2}R_x f(I(x)).
\]
Values at the origin are immaterial. This is the pullback of a one-form,
not componentwise composition with inversion. Direct differentiation gives
$DI(x)=|x|^{-2}R_x$, $R_x^2=\operatorname{Id}$,
$I(I(x))=x$, and $|\det DI(x)|=|x|^{-6}$.
Let $\mathcal G_3$ be the closed $L^3$ gradient space already used in the
manuscript, and let $w(f)$ be the unique minimum-norm representative of
$f+\mathcal G_3$, defined for all $f\in L^3$, solenoidal or not.

\begin{lemma}[Conformal covariance of the cubic representative]
\label{cm:covariance}
The map $\mathcal K$ is a linear isometric involution of $L^3(\R^3;\R^3)$,
it maps $\mathcal G_3$ onto itself, and
\begin{equation}\label{cm:commutation}
 w(\mathcal Kf)=\mathcal K w(f)\qquad(f\in L^3).
\end{equation}
\end{lemma}
\begin{proof}
Change variables $y=I(x)$; the cubic factor $|x|^{-6}$ cancels the
Jacobian, proving $\|\mathcal Kf\|_3^3=\|f\|_3^3$. The formulas above
also give $\mathcal K^2 f=f$ almost everywhere.

For $\phi\in C_c^\infty(\R^3)$ the chain rule away from zero gives
$\mathcal K\nabla\phi=\nabla(\phi\circ I)$. In fact $\phi\circ I$ is
identically zero near zero and extends smoothly there. Set
$\psi=\phi\circ I-\phi(0)$. At infinity, Taylor's theorem gives
$\psi(x)=O(|x|^{-1})$ and $\nabla\psi(x)=O(|x|^{-2})$. For a fixed smooth
cutoff $\chi$ equal to one on the unit ball and zero outside the ball of
radius two, the functions $\chi(x/R)\psi(x)$ are smooth and compactly
supported. Their gradients converge to $\nabla\psi$ in $L^3$:
the gradient tail tends to zero by integrability, while
\[
 \|R^{-1}(\nabla\chi)(\cdot/R)\psi\|_3
 \le C R^{-1}R^{-1}|\{R<|x|<2R\}|^{1/3}\le C/R.
\]
Thus $\mathcal K\nabla\phi\in\mathcal G_3$. Density and the isometry
extend this inclusion to $\mathcal G_3$; the involution makes it an
onto map. Consequently $\mathcal K$ maps the entire affine coset
$f+\mathcal G_3$ isometrically onto $\mathcal Kf+\mathcal G_3$.
Uniqueness of the cubic minimizer proves \eqref{cm:commutation}.
\end{proof}

\subsection{The weighted snapshot estimate}

We use the Hardy inequality
\begin{equation}\label{cm:hardy}
 \||x|^{-1}h\|_2\le2\|\nabla h\|_2
 \quad(h\in\dot H^1(\R^3)),
\end{equation}
componentwise for vector fields. Here the canonical homogeneous space has
its $L^6$ representative. To recall a proof, for a smooth compactly
supported scalar field, integrate $\operatorname{div}(x/|x|^2)=|x|^{-2}$
against $|h|^2$. Removal of a ball of radius $\varepsilon$ contributes
$O(\varepsilon)$; Cauchy--Schwarz then gives
$\int|h|^2/|x|^2\le2(\int|h|^2/|x|^2)^{1/2}\|\nabla h\|_2$.
Approximation in the gradient norm, and identification of the $L^6$
limit, proves \eqref{cm:hardy}. Summing the squared component inequalities
proves the vector version.

\begin{lemma}[Sobolev control after inversion]\label{cm:inverted-h1}
If $f\in H^1(\R^3;\R^3)$ and $|x|\nabla f\in L^2$, then
$\mathcal Kf\in H^1$ and, with $A=\||x|\nabla f\|_2$ and $U=\|f\|_2$,
\begin{equation}\label{cm:inverted-bound}
 \|\mathcal Kf\|_2\le2\|\nabla f\|_2,
 \qquad \|\nabla\mathcal Kf\|_2\le A+2\sqrt3\,U.
\end{equation}
\end{lemma}
\begin{proof}
Change of variables, initially for nonnegative integrands, gives
$\|\mathcal Kf\|_2^2=\int |y|^{-2}|f(y)|^2\,dy$, so the first bound is
\eqref{cm:hardy}. On every annulus the weak chain rule is legitimate and
\begin{equation}\label{cm:derivative}
 \nabla\mathcal Kf(x)
 =|x|^{-4}R_x(\nabla f)(I(x))R_x
     +|x|^{-3}\mathcal B(n_x,f(I(x))),
\end{equation}
where, for a unit vector $n$ and vector $v$,
\[
 \mathcal B(n,v)=-2v\otimes n-2(n\cdot v)\operatorname{Id}
                -2n\otimes v+8(n\cdot v)n\otimes n.
\]
Rotation to $n=e_1$ gives
\[
 |\mathcal B(n,v)|_F^2=8|v|^2+4(n\cdot v)^2\le12|v|^2.
\]
The first term on the right of \eqref{cm:derivative} has squared $L^2$
norm $\int|y|^2|\nabla f(y)|_F^2\,dy$, and the second has norm at most
$2\sqrt3\,U$. This proves the asserted bound for the punctured weak
gradient.

There is no distribution supported at the omitted point. Indeed, test
with a smooth function times a cutoff vanishing on $B_\varepsilon$ and
one outside $B_{2\varepsilon}$. The extra derivative pairing is bounded
by $C\|\mathcal Kf\|_{L^2(B_{2\varepsilon})}\varepsilon^{1/2}$, since
the $L^2$ norm of that cutoff's gradient is $O(\varepsilon^{1/2})$.
It tends to zero; all other terms pass by their local $L^2$
integrability. Thus the punctured derivatives are the global weak
derivatives, and $\mathcal Kf\in H^1$.
\end{proof}

\begin{theorem}[Finite energy and $L^{3/2}$ defect from a spatial moment]
\label{cm:snapshot}
Let $u\in H^1(\R^3;\R^3)$ be solenoidal and assume
$A=\||x|\nabla u\|_2<\infty$. Set $U=\|u\|_2$,
$c_0=\sqrt5/2$, $c_1=1+4\sqrt3$, and
$F=c_0(A+2\sqrt3 U)$. For its cubic representative $w=w(u)$,
$q=w-u$, and $\sigma=-\operatorname{div}w$, one has
\begin{align}
 \|w\|_2&\le2F,& \|q\|_2&\le2F+U,\label{cm:l2}\\
 \||x|\nabla w\|_2&\le c_1F,&
 \||x|\sigma\|_2&\le c_1F/\sqrt3.\label{cm:weighted}
\end{align}
In particular $w,q\in H^1$ and $\sigma\in L^{3/2}$, with
\begin{equation}\label{cm:sigma-threehalves}
 \|\sigma\|_{3/2}
 \le C_*\bigl(\|\sigma\|_2\,\||x|\sigma\|_2\bigr)^{1/2},
 \qquad C_*=2^{2/3}(4\pi/3)^{1/6}.
\end{equation}
The constants are not claimed sharp. No bound of $\|w\|_2$ by $U$
alone is asserted.
\end{theorem}
\begin{proof}
Although $\mathcal Ku$ need not be solenoidal, the Leray projection
$v=\mathbb P\mathcal Ku$ is in $H^1$, is solenoidal, and satisfies
$\|\nabla v\|_2\le\|\nabla\mathcal Ku\|_2$. Moreover
$(\operatorname{Id}-\mathbb P)\mathcal Ku\in\mathcal G_3$, by the
already established $L^3$ gradient decomposition. Thus
$w(v)=w(\mathcal Ku)=\mathcal Kw$.
Proposition~\ref{prop:quotient-divcurl}, applied to $v$, now gives
\begin{equation}\label{cm:inverted-representative}
 \mathcal Kw\in L^6,\qquad \|\nabla\mathcal Kw\|_2\le F.
\end{equation}
This use of the projection is essential: inversion itself does not
preserve ordinary solenoidality.

Apply \eqref{cm:hardy} to \eqref{cm:inverted-representative}. A second
change of variables gives the identity, valid also for infinite
nonnegative integrals,
\[
 \|w\|_2^2=\int\frac{|\mathcal Kw(x)|^2}{|x|^2}\,dx\le4F^2.
\]
This proves $w\in L^2$, rather than assuming it in a Hardy argument.
The estimate for $q$ follows by subtraction. The previously proved
unweighted div--curl estimate supplies $\nabla w,\nabla q\in L^2$.

Apply the annular chain rule \eqref{cm:derivative} to $w$ and rearrange it.
The same norm computation, using the just-proved $w\in L^2$, yields
\[
 \||x|\nabla w\|_2
 \le\|\nabla\mathcal Kw\|_2+2\sqrt3\|w\|_2\le c_1F.
\]
One can first integrate over annuli and let their endpoints tend to zero
and infinity; thus finiteness of the left side is a conclusion.
The nonlinear divergence constraint and Lemma~\ref{lem:trace-control}
give the pointwise inequality
$|\sigma|^2\le|\operatorname{sym}\nabla w|_F^2/3
\le|\nabla w|_F^2/3$. On $\{w=0\}$ its weak gradient vanishes almost
everywhere. Multiplication by $|x|^2$ proves the second bound of
\eqref{cm:weighted} without differentiating the unit direction.

For any $f$ with $f,|x|f\in L^2$ and any $R>0$, H\"older on the ball and
its complement gives
\[
 \int|f|^{3/2}\le(4\pi/3)^{1/4}
  \left(R^{3/4}\|f\|_2^{3/2}
       +R^{-3/4}\||x|f\|_2^{3/2}\right).
\]
Here $\int_{|x|>R}|x|^{-6}\,dx=(4\pi/3)R^{-3}$.
If either norm is zero then $f=0$ almost everywhere. Otherwise choose
$R=\||x|f\|_2/\|f\|_2$ and raise to power $2/3$.
This proves \eqref{cm:sigma-threehalves} with $f=\sigma$.
\end{proof}

\subsection{The weighted velocity budget on the original equation}

Fix the maximal branch of Proposition~\ref{prop:localtheory}. Put
\[
 E_0=\|u_0\|_2^2,\quad M_0=\||x|u_0\|_2^2,\quad
 Y(t)=\|\nabla u(t)\|_2^2,\quad
 M(t)=\||x|u(t)\|_2^2,\quad W(t)=\||x|\nabla u(t)\|_2^2.
\]
The last two integrals are initially allowed to be infinite. Let $S$ be
the homogeneous Sobolev constant used in Section~\ref{dc:section}.

\begin{proposition}[Input-only moment and weighted-gradient budget]
\label{cm:moment}
For every $H>0$, define
\begin{align*}
 K_H&=M_0+6\nu E_0H,\\
 I_H&=6S^{3/2}E_0^{1/4}H^{1/4}(E_0/(2\nu))^{3/4},\\
 L_H&=\tfrac12\left(I_H+\sqrt{I_H^2+4K_H}\right).
\end{align*}
Then for every $t<\min(H,T_*)$,
\begin{equation}\label{cm:moment-budget}
 M(t)+2\nu\int_0^tW(s)\,ds\le L_H^2.
\end{equation}
Moreover $W$ is bounded on every compact classical interval. Its latter
bound may depend on local high Sobolev norms; it is not claimed to be
uniform through $T_*$. The integrated estimate
\eqref{cm:moment-budget} is uniform through that endpoint and uses only
the displayed input quantities.
\end{proposition}
\begin{proof}
Use the bounded increasing weights
\[
 \rho_R(x)=\frac{|x|^2}{1+|x|^2/R^2},\qquad
 |\nabla\rho_R|\le2\sqrt{\rho_R},\qquad
 \Delta\rho_R=\frac{6-2|x|^2/R^2}{(1+|x|^2/R^2)^3}\le6.
\]
Set $M_R=\int\rho_R|u|^2$ and $W_R=\int\rho_R|\nabla u|^2$.
Testing the classical equation against $2\rho_Ru$, followed by integration
by parts, gives
\[
 M_R'+2\nu W_R
 =\nu\int\Delta\rho_R|u|^2
   +\int(\nabla\rho_R\cdot u)|u|^2
   +2\int p\,\nabla\rho_R\cdot u.
\]
For each fixed $R$ all coefficients are bounded and the local Sobolev
package justifies the identity by a further cutoff at infinity.
The pressure multiplier, as a linear map from matrix fields to scalars,
has Fourier symbol of Frobenius norm one. Consequently Plancherel gives
$\|p\|_2\le\|u\otimes u\|_2=\|u\|_4^2$.
Cauchy--Schwarz, interpolation and Sobolev therefore yield
\begin{equation}\label{cm:moment-differential}
 M_R'+2\nu W_R
 \le6\nu E_0+6\sqrt{M_R}\|u\|_4^2
 \le6\nu E_0+6S^{3/2}E_0^{1/4}\sqrt{M_R}Y^{3/4}.
\end{equation}
The factors in the middle expression are two from transport and four
from pressure. The energy identity gives
\[
 \int_0^tY^{3/4}\le H^{1/4}\left(\int_0^tY\right)^{3/4}
 \le H^{1/4}(E_0/(2\nu))^{3/4}.
\]
For a fixed compact time interval ending below $\min(H,T_*)$, let
$a_R$ be the supremum of $\sqrt{M_R}$ there. Integration of
\eqref{cm:moment-differential} gives $a_R^2\le K_H+I_Ha_R$.
Solving this quadratic yields $a_R\le L_H$ and, retaining dissipation,
$M_R(t)+2\nu\int_0^tW_R\le K_H+I_HL_H=L_H^2$.
Monotone convergence in $R$ proves \eqref{cm:moment-budget}.
No division by $E_0$ or $M_0$ is used, so zero data are included.

For pointwise weighted-gradient finiteness, differentiate the original
equation in each spatial direction and test against the corresponding
$2\rho_R\partial_k u$. The same cutoff justifications give
\[
 W_R'\le 2\|\nabla u\|_\infty W_R+6\nu Y
      +2\|u\|_\infty\sqrt{W_R}\sqrt Y
      +4\sqrt{W_R}\|\nabla p\|_2.
\]
The norm on the velocity gradient can be taken to be Frobenius.
The pressure multiplier also gives
$\|\nabla p\|_2\le2\|u\|_\infty\sqrt Y$.
For any fixed rate $\alpha>0$, Young's inequality therefore gives
\[
 W_R'\le(2\|\nabla u\|_\infty+\alpha)W_R
       +(6\nu+25\|u\|_\infty^2/\alpha)Y.
\]
All coefficients are bounded on each compact classical interval, and
$W_R(0)\le\||x|\nabla u_0\|_2^2<\infty$. Gronwall, followed by monotone
convergence in $R$, proves the stated local bound for $W$.
This step establishes pointwise applicability of Theorem~\ref{cm:snapshot};
it is not used as an endpoint-uniform producer.
\end{proof}

\begin{corollary}[Unconditional moment-based spacetime bounds]
\label{cm:spacetime}
For every $H>0$ and every $t<\min(H,T_*)$, set
\[
 T_H=\frac{L_H^2}{2\nu}+12HE_0,
 \qquad B_H=\frac56(1+4\sqrt3)^2T_H.
\]
Then
\begin{align}
 \int_0^t\|w(s)\|_2^2\,ds&\le10T_H,\label{cm:w-budget}\\
 \int_0^t\||x|\sigma(s)\|_2^2\,ds&\le B_H,\label{cm:sigma-weight-budget}\\
 \int_0^t\|\sigma(s)\|_{3/2}^2\,ds
 &\le C_*^2\sqrt{\frac{E_0B_H}{8\nu}}.\label{cm:sigma-spacetime}
\end{align}
Thus $w(t)\in L^2$ and $\sigma(t)\in L^{3/2}$ at every classical time,
and the displayed time integrals have input-only bounds through any
putative finite endpoint. Also $q=w-u\in L^2_tL^2_x$ on those intervals.
\end{corollary}
\begin{proof}
Theorem~\ref{cm:snapshot} applies at every time by
Proposition~\ref{cm:moment}. Since
$(A+2\sqrt3 U)^2\le2A^2+24U^2$ and $c_0^2=5/4$, its estimates give
\[
 \|w\|_2^2\le10(W+12E),\qquad
 \||x|\sigma\|_2^2\le\tfrac56(1+4\sqrt3)^2(W+12E).
\]
Integrate, use $E\le E_0$ and \eqref{cm:moment-budget}, and apply
Cauchy--Schwarz in time to the square of
\eqref{cm:sigma-threehalves}. The already proved defect budget is
$\int_0^t\|\sigma\|_2^2\le E_0/(8\nu)$, giving
\eqref{cm:sigma-spacetime}. The assertion for $q$ follows from
$\|q\|_2^2\le2\|w\|_2^2+2E$.

For measurability, $w(t)$ is continuous into $L^3$ by the existing
minimizer stability theorem. Its distributional derivatives are weakly
measurable by pairing with smooth compactly supported tests. On each
ball, their weighted norms are supremums over countable dense sets of
dual tests; exhaustion of space gives Borel measurability of all extended
norms used above. The same argument applies to the $L^2$ norm of $w$ and
the $L^{3/2}$ norm of $\sigma$. Thus the integrals are well defined.
\end{proof}

\subsection{Why this does not close the temporal producer}

The critical defect line is $2/s+3/a=2$, because under Navier--Stokes
scaling $\sigma_\lambda(t,x)=\lambda^2\sigma(\lambda^2t,\lambda x)$.
The new bound has $(s,a)=(2,3/2)$ and left side $3$, not $2$.
In the usual Navier--Stokes terminology this is a supercritical,
not a subcritical, bound. It gives neither a uniform small $L^{3/2}$
defect nor the critical
fourth-power time integral of its $L^2$ norm. This is not a defect of the
endpoint justification: the time integrals themselves are uniform, but
are of the wrong scaling strength.

\begin{remark}[A self-similar curve testing only the proposed inference]
\label{cm:scaling-test}
The following is an explicit curve of fields, not a Navier--Stokes
solution. It tests the implication from the new integrability budgets
alone, not the signed vector equation. Let
\[
 U(x)=(4x_2,-2x_1,0)\exp[-(x_1^2+2x_2^2+x_3^2)].
\]
It is Schwartz and solenoidal, and
$\operatorname{div}(|U|U)(1,1,0)=-8e^{-6}/\sqrt5\ne0$.
Its minimizing defect $\sigma_U$ is nonzero: otherwise $w(U)$ would
be both solenoidal and in $U+\mathcal G_3$, hence equal to $U$ by the
Leray decomposition, contradicting the variational divergence condition.
Let $E_U=\|U\|_2^2$, $Y_U=\|\nabla U\|_2^2$, fix $T>0$, and define
\[
 \lambda(t)=\left(4\nu Y_U(T-t)/E_U\right)^{-1/2},\qquad
 v(t,x)=\lambda(t)U(\lambda(t)x),\qquad 0\le t<T.
\]
The homogeneity and dilation covariance of the cubic minimizer imply
$w(v)=\lambda w(U)(\lambda x)$ and
$\sigma_v=\lambda^2\sigma_U(\lambda x)$. In particular
\[
 \|\sigma_v\|_2^2=\lambda\|\sigma_U\|_2^2,\quad
 \||x|\sigma_v\|_2^2=\lambda^{-1}\||x|\sigma_U\|_2^2,\quad
 \|\sigma_v\|_{3/2}=\|\sigma_U\|_{3/2}.
\]
All three squared time integrals in the new and old defect budgets are
finite; so is $\int_0^T\|w(v)\|_2^2$. But
$\int_0^T\|\sigma_v\|_2^4=\infty$. Moreover the choice of $\lambda$
makes the ordinary energy identity exact:
\[
 \frac{d}{dt}\|v\|_2^2+2\nu\|\nabla v\|_2^2
 =-E_U\lambda'/\lambda^2+2\nu Y_U\lambda=0.
\]
The spatial moment and weighted-gradient integrals are finite as well,
since they scale respectively as $\lambda^{-3}$ and $\lambda^{-1}$.
This example asserts finiteness of the budgets, not every numerical
constant in \eqref{cm:moment-budget} for the curve. It is a backward
self-similar test curve of the same type as the earlier research
countermodels, not a new singular solution or a novelty claim. Its
constant critical velocity norm already prevents it from being a
singular selected branch by the manuscript's endpoint theorem.
\end{remark}

\paragraph{Remaining obligation.}
The spatial questions are answered here with weighted data, and the
weights have an unconditional integrated budget on the original equation.
The arbitrary-data critical temporal estimate remains unproved. A proof
must still use additional trajectory information to obtain signed
pressure/strain control or a critical defect budget; it cannot replace
that step by the new spatial membership assertions.
